% -------------------------------------------------------------------------
% WBS ID: RES-NUM-WET
% Deliverable: Wetting–Drying, Positivity and Discontinuity Treatment
% Status: Draft complete
% Evidence status: Regional research specification complete
% Review status: Not reviewed
% Last updated: 2026-07-07
% Dependencies: RES-PHY-NRS, RES-NUM-FLX, RES-NUM-SRC, RES-NUM-TIM
% Downstream interfaces: RES-VER-BMK, RES-VER-TST, Software wet/dry implementation
% -------------------------------------------------------------------------

\section{RES-NUM-WET --- Wetting–Drying, Positivity and Discontinuity Treatment}
\label{sec:res-num-wet}

\subsection{Purpose and Scope}
\label{sec:res-num-wet-purpose}

% TODO: Define what this Deliverable covers and explicitly excludes.

\subsection{Research Questions}
\label{sec:res-num-wet-research-questions}

% TODO: State the questions that must be answered before the Deliverable
% can be accepted.

\subsection{Terminology and Definitions}
\label{sec:res-num-wet-definitions}

% TODO: Define the technical terminology, symbols, quantities and
% classifications used in this Deliverable.

\subsection{Literature Evidence}
\label{sec:res-num-wet-literature-evidence}

% TODO: Record evidence from peer-reviewed papers, authoritative datasets,
% standards, technical reports and primary documentation.
%
% Do not create a detached literature-summary list. Explain how each source
% supports, contradicts or limits a project decision.

\subsection{Governing Relations and Quantitative Definitions}
\label{sec:res-num-wet-governing-relations}

% TODO: Record relevant equations, dimensionless groups, empirical models,
% extraction rules or quantitative definitions.
%
% Use align environments for displayed equations.
% Every displayed equation must have a unique \label{eq:...}.
% Do not place punctuation after displayed equations.

\subsection{Discrete Formulation}
\label{sec:res-num-wet-discrete-formulation}

The regional model uses positivity-preserving wetting--drying logic and
does not use regional VOF. Positivity preservation enforces
$h_i\ge 0$; wetting--drying additionally controls shoreline motion,
near-dry velocities, reconstruction, flux fallback, draining and dry
reset.

The selected dry reset and draining limiter are
\cref{eq:res-num-spec-regional-dry-reset,eq:res-num-spec-regional-draining-limiter}.

\subsection{Conservation Properties}
\label{sec:res-num-wet-conservation-properties}

Outgoing donor fluxes shall be scaled consistently so the same limited
face flux is applied to both adjacent cells with opposite signs. Dry
resets shall be audited as controlled positivity corrections rather than
silent conservation-preserving fluxes.

\subsection{Consistency and Formal Accuracy}
\label{sec:res-num-wet-consistency-formal-accuracy}

% TODO: Complete this RES-NUM Domain-specific subsection.

\subsection{Stability Requirements}
\label{sec:res-num-wet-stability-requirements}

% TODO: Complete this RES-NUM Domain-specific subsection.

\subsection{Mesh and Time-Step Requirements}
\label{sec:res-num-wet-mesh-time-step-requirements}

% TODO: Complete this RES-NUM Domain-specific subsection.

\subsection{Numerical Failure Modes}
\label{sec:res-num-wet-numerical-failure-modes}

% TODO: Complete this RES-NUM Domain-specific subsection.

\subsection{Verification Requirements}
\label{sec:res-num-wet-verification-requirements}

% TODO: Complete this RES-NUM Domain-specific subsection.

\subsection{Candidate Approaches}
\label{sec:res-num-wet-candidate-approaches}

% TODO: Define the credible candidate approaches considered.

\subsection{Critical Comparison}
\label{sec:res-num-wet-critical-comparison}

% TODO: Compare candidates using physically and project-relevant criteria.
% Do not select an approach solely on popularity or implementation ease.

\subsection{Assumptions and Applicability}
\label{sec:res-num-wet-assumptions}

% TODO: State assumptions and the conditions under which they are valid.

\subsection{Limitations and Failure Conditions}
\label{sec:res-num-wet-limitations}

% TODO: State where the evidence, model or selected approach ceases to be
% reliable or applicable.

\subsection{Project-Specific Conclusions}
\label{sec:res-num-wet-conclusions}

% TODO: Convert the evidence into conclusions specific to the Studentship
% project. Do not merely restate literature findings.

\subsection{Selected Approach and Rationale}
\label{sec:res-num-wet-selected-approach}

% TODO: Record the selected approach only after sufficient evidence exists.
% State the alternatives rejected and the reasons for rejection.

The selected regional wet/dry treatment combines hydrostatic
reconstruction, HLL/HLLE fallback, first-order near-dry reconstruction,
a draining limiter, velocity regularisation and post-stage dry reset.

\subsection{Unresolved Decisions}
\label{sec:res-num-wet-unresolved-decisions}

% TODO: Record unresolved questions, required evidence and decision owners.

Final values of $h_{\mathrm{dry}}$, velocity-regularisation depth,
near-dry reconstruction threshold and mass-correction audit tolerances
remain verification parameters.

\subsection{Requirements and Handoffs}
\label{sec:res-num-wet-requirements}

% TODO: State requirements passed to other Research Domains, Software,
% verification activities or later execution work.

\subsection{Deliverable Completion Record}
\label{sec:res-num-wet-completion-record}

% TODO: Record whether the Deliverable is:
%
% - Not started
% - In progress
% - Draft complete
% - Reviewed
% - Accepted
%
% Acceptance requires:
%
% - the research questions to be answered;
% - claims to be supported by appropriate evidence;
% - assumptions and limitations to be explicit;
% - the project-specific conclusion to be stated;
% - downstream requirements to be traceable;
% - unresolved decisions to be closed or formally deferred.
